% GENERATED FILE. Edit canonical lessons, then run npm run generate:books.
% canonical-source-hash: fnv1a64:f05fca59ea799c63
% canonical-lessons: BN-C502-kompiutar, BN-C502-mobail, BN-C502-selphi, BN-C502-ghata, BN-C502-charj-deoya

\chapter{Screens and Phones}
\label{ch:bn-screens-and-phones}
\hlchaptermodality{502}

\begin{chapteropening}
\textbf{By the end of this chapter:} \emph{I can talk about computers, phones and videos.}

Complete the last lesson of chapter 502: I can talk about computers, phones and videos.
\end{chapteropening}

\section[\texorpdfstring{\bn{কম্পিউটার} (kômpiuṭār)}{কম্পিউটার (kômpiuṭār)}]{\bn{কম্পিউটার} (kômpiuṭār) — a computer}

\label{lesson:BN-C502-kompiutar}



\begin{quote}
Before the new one: say the Bengali for an architect, then the Bengali for to diagnose.
\end{quote}

\subsection*{You'll want to know: \bn{কম্পিউটার}}
\textbf{\bn{কম্পিউটার}} — \emph{kômpiuṭār} — \textquotedblleft{}a computer\textquotedblright{}.

\begin{quote}
\bn{আমি} \bn{অফিসে} \bn{সারাদিন} \bn{কম্পিউটারে} \bn{কাজ} \bn{করি।}
\end{quote}

\begin{quote}
\emph{Āmi ôphise sārādin kômpiuṭāre kāj kôri.}
\end{quote}

\begin{quote}
I work at a computer all day at the office.
\end{quote}

\begin{cousinweb}
\textbf{\bn{কম্পিউটার}} (\emph{kômpiuṭār}) is borrowed from English \emph{computer}.
\end{cousinweb}

\begin{grammarlens}[title={What you've built}]
One more word for describing what you lived and felt.
\end{grammarlens}

\subsection*{Your turn}
\begin{itemize}
\raggedright
  \item \emph{Say it:} \emph{kômpiuṭār}
  \item \emph{Say it:} \emph{kômpiuṭār}, once more
  \item \emph{Say it:} \emph{Āmi ôphise sārādin kômpiuṭāre kāj kôri.} — I work at a computer all day at the office.
  \item \emph{Recall:} say the Bengali for an architect, then the Bengali for to diagnose, then say \emph{kômpiuṭār} again
\end{itemize}

\begin{tcolorbox}[breakable,colback=teal!4,colframe=teal!35!black,title={Before you move on}]
What does \bn{কম্পিউটার} mean? (\textquotedblleft{}A computer\textquotedblright{}.) Say it once more.
\end{tcolorbox}
\section[\texorpdfstring{\bn{মোবাইল} (mobāil)}{মোবাইল (mobāil)}]{\bn{মোবাইল} (mobāil) — a mobile phone}

\label{lesson:BN-C502-mobail}



\begin{quote}
Before the new one: say the Bengali for to diagnose, then the Bengali for a computer.
\end{quote}

\subsection*{You'll want to know: \bn{মোবাইল}}
\textbf{\bn{মোবাইল}} — \emph{mobāil} — \textquotedblleft{}a mobile phone\textquotedblright{}.

\begin{quote}
\bn{আমার} \bn{মোবাইলের} \bn{ব্যাটারি} \bn{খারাপ।}
\end{quote}

\begin{quote}
\emph{Āmār mobāiler bæṭāri khārāp.}
\end{quote}

\begin{quote}
My mobile's battery is bad.
\end{quote}

\begin{cousinweb}
\textbf{\bn{মোবাইল}} (\emph{mobāil}) is borrowed from English \emph{mobile}, short for a mobile phone.
\end{cousinweb}

\begin{grammarlens}[title={What you've built}]
One more word for describing what you lived and felt.
\end{grammarlens}

\subsection*{Your turn}
\begin{itemize}
\raggedright
  \item \emph{Say it:} \emph{mobāil}
  \item \emph{Say it:} \emph{mobāil}, once more
  \item \emph{Say it:} \emph{Āmār mobāiler bæṭāri khārāp.} — My mobile's battery is bad.
  \item \emph{Recall:} say the Bengali for to diagnose, then the Bengali for a computer, then say \emph{mobāil} again
\end{itemize}

\begin{tcolorbox}[breakable,colback=teal!4,colframe=teal!35!black,title={Before you move on}]
What does \bn{মোবাইল} mean? (\textquotedblleft{}A mobile phone\textquotedblright{}.) Say it once more.
\end{tcolorbox}
\section[\texorpdfstring{\bn{সেলফি} (selphi)}{সেলফি (selphi)}]{\bn{সেলফি} (selphi) — a selfie}

\label{lesson:BN-C502-selphi}



\begin{quote}
Before the new one: say the Bengali for a computer, then the Bengali for a mobile phone.
\end{quote}

\subsection*{You'll want to know: \bn{সেলফি}}
\textbf{\bn{সেলফি}} — \emph{selphi} — \textquotedblleft{}a selfie\textquotedblright{}.

\begin{quote}
\bn{মেয়েটা} \bn{ফোনে} \bn{একটা} \bn{সেলফি} \bn{তুলল।}
\end{quote}

\begin{quote}
\emph{Meyeṭā phone ekṭā selphi tullô.}
\end{quote}

\begin{quote}
The girl took a selfie with her phone.
\end{quote}

\begin{cousinweb}
\textbf{\bn{সেলফি}} (\emph{selphi}) is borrowed from English \emph{selfie}.
\end{cousinweb}

\begin{grammarlens}[title={What you've built}]
One more word for describing what you lived and felt.
\end{grammarlens}

\subsection*{Your turn}
\begin{itemize}
\raggedright
  \item \emph{Say it:} \emph{selphi}
  \item \emph{Say it:} \emph{selphi}, once more
  \item \emph{Say it:} \emph{Meyeṭā phone ekṭā selphi tullô.} — The girl took a selfie with her phone.
  \item \emph{Recall:} say the Bengali for a computer, then the Bengali for a mobile phone, then say \emph{selphi} again
\end{itemize}

\begin{tcolorbox}[breakable,colback=teal!4,colframe=teal!35!black,title={Before you move on}]
What does \bn{সেলফি} mean? (\textquotedblleft{}A selfie\textquotedblright{}.) Say it once more.
\end{tcolorbox}
\section[\texorpdfstring{\bn{ঘাঁটা} (ghā̃ṭā)}{ঘাঁটা (ghā̃ṭā)}]{\bn{ঘাঁটা} (ghā̃ṭā) — to rummage; to fiddle with}

\label{lesson:BN-C502-ghata}



\begin{quote}
Before the new one: say the Bengali for a mobile phone, then the Bengali for a selfie.
\end{quote}

\subsection*{You'll want to know: \bn{ঘাঁটা}}
\textbf{\bn{ঘাঁটা}} — \emph{ghā̃ṭā} — \textquotedblleft{}to rummage; to fiddle with\textquotedblright{}.

\begin{quote}
\bn{ছেলেটা} \bn{সারাদিন} \bn{ফোন} \bn{ঘাঁটে।}
\end{quote}

\begin{quote}
\emph{Chheleṭā sārādin phon ghā̃ṭe.}
\end{quote}

\begin{quote}
The boy fiddles with his phone all day.
\end{quote}

Past (he/she …): \bn{তিনি} \bn{আলমারি} \bn{ঘাঁটলেন}, \emph{tini ālmāri ghā̃ṭlen}.

\begin{cousinweb}
\textbf{\bn{ঘাঁটা}} (\emph{ghā̃ṭā}) is a tadbhava: CDIAL lists Bengali \emph{ghā̃ṭā}, \textquotedblleft{}to stir with the fingers\textquotedblright{}, under Sanskrit \emph{ghaṭṭayati} (CDIAL 4417).
\end{cousinweb}

\begin{grammarlens}[title={What you've built}]
One more word for describing what you lived and felt.
\end{grammarlens}

\subsection*{Your turn}
\begin{itemize}
\raggedright
  \item \emph{Say it:} \emph{ghā̃ṭā}
  \item \emph{Say it:} \emph{ghā̃ṭā}, once more
  \item \emph{Say it:} \emph{Chheleṭā sārādin phon ghā̃ṭe.} — The boy fiddles with his phone all day.
  \item \emph{Recall:} say the Bengali for a mobile phone, then the Bengali for a selfie, then say \emph{ghā̃ṭā} again
\end{itemize}

\begin{tcolorbox}[breakable,colback=teal!4,colframe=teal!35!black,title={Before you move on}]
What does \bn{ঘাঁটা} mean? (\textquotedblleft{}To rummage; to fiddle with\textquotedblright{}.) Say it once more.
\end{tcolorbox}
\section[\texorpdfstring{\bn{চার্জ} \bn{দেওয়া} (chārj deoyā)}{চার্জ দেওয়া (chārj deoyā)}]{\bn{চার্জ} \bn{দেওয়া} (chārj deoyā) — to charge a phone; to charge (a battery)}

\label{lesson:BN-C502-charj-deoya}



\begin{quote}
Before the new one: say the Bengali for a selfie, then the Bengali for to rummage.
\end{quote}

\subsection*{You'll want to know: \bn{চার্জ} \bn{দেওয়া}}
\textbf{\bn{চার্জ} \bn{দেওয়া}} — \emph{chārj deoyā} — \textquotedblleft{}to charge a phone; to charge (a battery)\textquotedblright{}.

\begin{quote}
\bn{রাতে} \bn{আমি} \bn{ফোনে} \bn{চার্জ} \bn{দিলাম।}
\end{quote}

\begin{quote}
\emph{Rāte āmi phone chārj dilām.}
\end{quote}

\begin{quote}
I charged my phone at night.
\end{quote}

Past (he/she …): \bn{তিনি} \bn{ল্যাপটপে} \bn{চার্জ} \bn{দিলেন}, \emph{tini læpṭôpe chārj dilen}.

\begin{cousinweb}
\textbf{\bn{চার্জ} \bn{দেওয়া}} (\emph{chārj deoyā}) joins \emph{charge}, borrowed from English, to \bn{দেওয়া}, \textquotedblleft{}to give\textquotedblright{}: giving the battery its charge.
\end{cousinweb}

\begin{grammarlens}[title={What you've built}]
One more word for describing what you lived and felt.
\end{grammarlens}

\subsection*{Your turn}
\begin{itemize}
\raggedright
  \item \emph{Say it:} \emph{chārj deoyā}
  \item \emph{Say it:} \emph{chārj deoyā}, once more
  \item \emph{Say it:} \emph{Rāte āmi phone chārj dilām.} — I charged my phone at night.
  \item \emph{Recall:} say the Bengali for a selfie, then the Bengali for to rummage, then say \emph{chārj deoyā} again
\end{itemize}

\begin{tcolorbox}[breakable,colback=teal!4,colframe=teal!35!black,title={Before you move on}]
What does \bn{চার্জ} \bn{দেওয়া} mean? (\textquotedblleft{}To charge a phone; to charge (a battery)\textquotedblright{}.) Say it once more.
\end{tcolorbox}
